\section{Openness above the marginal coefficient}\label{sec:openness}

We now exclude a supermarginal coefficient at the mass threshold.
The required assertion concerns the original Skellam height law:
\begin{proposition}\label{prop:openness}
If \(a(b)>8\pi\), there is an open interval \(I\) containing \(b\),
contained in \((0,\infty)\), such that \(a(\widetilde b)>0\) for
every \(\widetilde b\in I\).
\end{proposition}
We will not assume continuity of \(a\), nor differentiate it.  The
argument first adds a positive precision on \emph{block averages}.
This supplies a Gaussian large-field bound for the observed field,
while its scaling coefficient can be kept above \(8\pi\).
We prove entry into the small-activity maps at fixed reference, and
then use a torus variance diagnostic to return to the original model.

\subsection{A strengthened law and its local limits}

Fix \(b\) with \(a=a(b)>8\pi\).  Choose \(\delta>0\) sufficiently small
that
\begin{equation}\label{eq:op-reference}
 a'=(a^{-1}+\delta)^{-1}>8\pi,\qquad \alpha=\sqrt{a'}.
\end{equation}
Both \(\delta\) and \(a'\) remain fixed when the physical parameter
\(\widetilde b\) varies.  For a dyadic \(s_0\), let \(B_{s_0}\) average
over disjoint square cells of side \(s_0\), and let \(A_{\rm cell}\)
be the unit-conductance nearest-neighbor Laplacian on these cells.
Add to the height energy the term
\begin{equation}\label{eq:op-added-precision}
 \frac{\delta}{2}\ip{B_{s_0}h}{A_{\rm cell}B_{s_0}h}.
\end{equation}
All grids are nested and have the common square-symmetry origins
used below.  Set \(s=ts_0\), where \(t\) is also dyadic, and observe
\begin{equation}\label{eq:op-observation}
 Y=B_sh+\eta,\qquad \eta\sim N(0,\Id),
\end{equation}
with independent observation noise.  On tori we work modulo global
height translation by \(2\pi\).  When local observations confine a
component constant, its translations are summed instead.

The comparison observation is \(B_s\alpha\psi+\eta\), where the
mean-free covariance of \(\psi\) is \(A^+\) and its constant is
uniform modulo \(2\pi/\alpha\).  In particular, this is the
\emph{ordinary nearest-neighbor} Gaussian reference; no averaged
precision is added to it.  On a cut graph \(G\), write its minimized
observation energy as
\begin{equation}\label{eq:op-observation-energy}
 E_G(y)=\inf_v\left\{a'^{-1}\ip{v}{A_Gv}
                         +\norm{B_sv-y}_2^2\right\},
\end{equation}
using free constants on the components.  Denote by \(Z_G(y;\widetilde b)\)
the corresponding Skellam partition sum with
\eqref{eq:op-added-precision} and observation penalty
\(\norm{B_sh-y}_2^2/2\).

The local domains are the mixed polygonal patches of
Theorem~\ref{thm:mixed-limits}, aligned with whole \(s\)-cells.
Whenever microscopic bonds between tiles are deleted, delete also
the bonds of \(A_{\rm cell}\) crossing that cut.  Thus a pre-pin
component has precisely its own subcells and face bonds.  A hard
core pin used to dominate a restoration includes all primary
vertices in every subcell supporting a restored averaged bond.
Such supports fit inside the fixed core collars.

\begin{lemma}[Strengthened local limits]\label{lem:op-local-limits}
Fix a bounded mixed geometry in observation-cell units, with
separated positive-area core and guard pins and observation masses
controlling each unpinned constant.  At the central parameter \(b\),
the strengthened laws have the Gaussian smear, conditional-smear,
and pin-interaction limits of Theorems~\ref{thm:mixed-limits}
and~\ref{thm:pin-interaction}, with coefficient \(a'\), in the
iterated limit
\[
             \lim_{t\to\infty}\lim_{s_0\to\infty}.
\]
The limits hold under each indicated pin set.  The quantities
\(D\) and \(r_H\) retain their centered-law definitions;
finite-probe convergence of these conditional quantities holds
also when their data are sampled under any fixed real observation
tilt.  This is convergence in probability under the corresponding
guard laws, not a redefinition of the centered pin ratios.
At fixed geometry the conclusions are uniform on compact sets of
observation data.
\end{lemma}

\begin{proof}
For fixed \(t\), the additional energy is a finite-rank quadratic
tilt of finitely many subcell averages.  Theorems~\ref{thm:mixed-limits}
and~\ref{thm:pin-interaction} give its limit on each fixed pin set.
Integration against the nonnegative quadratic weight is justified
by its boundedness and the confining observation masses; fixed
linear tilts are handled by the exponential moment bounds.
Free component constants are treated by the joint uniform-phase
limit before summing their confining translations.

The full pin interaction also transforms by finite-dimensional
quantities.  If \(W\) is the added quadratic weight, and subscripts
indicate centered conditioning on the designated pins, then
\begin{equation}\label{eq:op-tilted-interaction}
 D_W(H,p)-D(H,p)
 =\log\E_{H,p}W+\log\E W
       -\log\E_HW-\log\E_pW .
\end{equation}
Each expectation has a strictly positive Gaussian limit.  This
proves, at fixed \(t\), convergence to the pin interaction for the
tilted continuum Gaussian precision, not merely convergence of
unconditioned smears.

In observation-cell units, let \(P_t\) denote averages over squares
of side \(1/t\), and \(B\) averages over the unit observation cells.
The continuum form obtained in this first limit is
\begin{equation}\label{eq:op-continuum-form}
 T_t(f)=a^{-1}\int|\nabla f|^2
       +\delta\sum_{\langle i,j\rangle}
                  \bigl((P_tf)_i-(P_tf)_j\bigr)^2
       +\norm{Bf}_2^2 .
\end{equation}
It converges in inverse forms to
\begin{equation}\label{eq:op-limit-form}
                   T(f)=a'^{-1}\int|\nabla f|^2+\norm{Bf}_2^2,
\end{equation}
also with any of the specified zero-region pins.  Here is the
form argument, including its normalization.  If \(\zeta=1/t\),
the averaged-gradient term is a sum of
\(\zeta^2|D_\zeta P_tf|^2\).  Bounded \(T_t\)-energy gives weak
\(H^1\) and strong \(L^2\) compactness on each fixed component.
The averaged difference quotients converge weakly to the
directional derivatives on its interior, giving the lower bound
with coefficient \(\delta\).  Conversely, smooth functions give
the matching upper bound by Riemann sums.  The averaged-gradient
term is uniformly bounded by a constant times continuum gradient
energy, by the two-cell Poincare inequality.  Energy approximation
therefore extends recovery to all admissible functions, using
corner cutoffs and density with zero neighborhoods of the pins,
as in \Bcite{an:forms}.  No component is glued through a point.
This proves the claimed inverse-form convergence and, in
particular, convergence of all finite covariance and pinned
inverse matrices.

We also need convergence of the infinite-dimensional interaction
determinant.  Form convergence alone would not give it.  We verify
the uniform compactness used in \Bcite{an:gaussian-localization}
for the additional averaged term.
For an energy-harmonic function off a pin or source, test its
equation against \(\chi^2f\), where \(\chi\) is a smooth cutoff
in a separating margin.  On a subcell replace
\((P_t(\chi^2f))_i\) by \(\chi_i^2(P_tf)_i\).  The error satisfies
\begin{equation}\label{eq:op-cutoff-error}
 \left|(P_t(\chi^2f))_i-\chi_i^2(P_tf)_i\right|
 \leq C\bigl(t^{-1}|\chi_i|+t^{-2}\bigr)
                         (P_tf^2)_i^{1/2}.
\end{equation}
Pairing this error with each incident averaged bond and applying
Cauchy--Schwarz absorbs the term with \(|\chi_i|\) into a fixed
part of the positive weighted bond energy.  The remaining
\(t^{-2}\) terms cost local \(L^2\) mass: use
\[
 |(P_tf)_i-(P_tf)_j|
       \leq(P_tf^2)_i^{1/2}+(P_tf^2)_j^{1/2}
\]
and the bounded number of incident bonds.  The differences of
the cutoff values themselves are controlled by the discrete
product inequality.  All bounds are uniform for large \(t\).

The unit-cell observation term may cross an intermediate cutoff.
On the fixed geometry its forcing consists of finitely many
bounded cell-average coefficients.  Carry these coefficients
as finite-dimensional coordinates through the successive
restriction maps.  Their bounds follow from the original energy
bound.  The remaining terms are local on an enlarged margin for
large \(t\).  We consequently obtain the same nested
gradient-from-\(L^2\) estimates as for the form with unit-cell
mass alone, with this fixed finite-rank enlargement.

Two broken-domain \(H^1\)-to-\(L^2\) compact embeddings, each
with singular numbers \(O(j^{-1/2})\), separated by these
harmonic interior estimates, give a uniform \(O(j^{-1})\)
singular-value tail for the cross projections between the
separated pin regions.  A fixed finite-rank enlargement changes
only finitely many singular directions.  Bounded-energy
sequences harmonic off one pin have weak limits harmonic off
that pin for \(T\): test their equations against smooth
functions, whose averaged differences converge strongly.
For a weakly null sequence, compact \(L^2\) convergence and the
interior estimate give strongly vanishing energy near the
other pin; the observation averages vanish as well.

For completeness, identify all the varying energy spaces with
the same broken \(H^1\) space \(V\).  The forms \(T_t\) are
uniformly equivalent there, since the original gradient and
observation terms are a lower bound and the averaged-gradient
term has a uniform upper bound.  Let \(V_{0,H}\) be the fixed
subspace of functions zero on \(H\).  Its \(T_t\)-orthogonal
complement \(U_H^t\) is the source projection space.  If
\(\phi\) is a smooth density supported in \(H\), orthogonality
of \(v_t\in U_H^t\) to the \(T_t\)-Riesz representative of
\(\phi\) is exactly \(\int v_t\phi=0\), independently of \(t\).

Take \(T_t\)-unit vectors in \(U_H^t\) orthogonal to increasing
finite lists of these probes.  Every weak limit vanishes on
\(H\) and is \(T\)-harmonic off \(H\), by the smooth-test
passage above.  It lies both in \(V_{0,H}\) and its
\(T\)-orthogonal complement, so is zero.  The finite
observation forcing coordinates are the actual averages
\(Bv_t\); strong \(L^2\) convergence makes them zero too.
For a cutoff \(\chi_p\) equal to one near the opposite pin,
the local estimate gives \(T_t(\chi_pv_t)\to0\).
Moreover
\(\norm{P_p^tv_t}_{T_t}\leq\norm{\chi_pv_t}_{T_t}\),
since \(v_t-\chi_pv_t\) vanishes on \(p\).  This proves
uniform finite-probe approximation of the cross projection.
The same argument applies in reverse, and the finite-probe
Gram matrices converge by inverse-form convergence.
Combined with the uniform square-summable singular tails,
these facts give convergence of the cross determinants.
The limiting cross projection has norm strictly smaller than
one, as in \Bcite{an:gaussian-localization}; operator convergence
preserves a norm gap.  In particular the logarithmic determinant
defining \(D(H,p)\) converges.

The conditional conclusions now follow by the regression and
minimal-Hessian argument used in Theorem~\ref{thm:pin-interaction},
with the covariance and determinant just obtained.  Conditional
Laplace transforms are first approximated using finitely many
smooth guard probes.  The full pin ratio follows by increasing
finite core penalties and using the centered identity
\(\E r_H=e^{-D(H,p)}\).
For use in both neighboring restoration topologies, take a
common finite probe list.  The bound
\(r_H\leq r_\Delta\leq1\) from Lemma~\ref{lem:restoration}, and
strict positivity of the Gaussian limiting \(r_H\), transfer
negligible exceptional sets in both directions.  Fixed real
tilts preserve convergence for these same centered-defined
quantities by the exponential moment bounds; the centered
expectation identity is not asserted under tilted sampling.

Every geometry and probe list is fixed before refinement.
Compact-data uniformity follows by taking convergent sequences
of observation vectors and using the same exponential
integrability.  Operationally, given an error tolerance one
first chooses a sufficiently large \(t\), then a sufficiently
large \(s_0\) depending on \(t\).  No single refinement threshold
uniform over all \(t\) is required.
\end{proof}

\subsection{A refinement-independent large-field reserve}

The bare quadratic energy of the subdivided Skellam law cannot
be used as a uniform site-level Gaussian tail.  The next bound
is the reason for adding \eqref{eq:op-added-precision}.

\begin{lemma}\label{lem:op-large-field}
For every cut topology used in the entry construction, there is
\(c_*>0\), depending only on \(\delta,a'\) and the fixed local
geometry conventions, such that
\begin{equation}\label{eq:op-large-field}
 \frac{Z_G(y;\widetilde b)}{Z_G(0;\widetilde b)}
                  \leq \exp\{-c_*E_G(y)/2\}.
\end{equation}
The same constant works for all sufficiently large dyadic
refinements and for \(\widetilde b\) in a neighborhood of \(b\).
\end{lemma}

\begin{proof}
On the subdivided quadratic lattice, complete the square in the
joint height and observation variables.  Centered theta
domination bounds the shifted sum by the centered sum times
the exponential of minus half its real minimum energy.
Discard the nonnegative microscopic chain energy in that
minimum, retaining the added average energy and
\(\norm{B_sh-y}_2^2\).  We show that this retained energy
dominates \(c_*E_G(y)\) for every real trial height.

Write \(u_i=(B_{s_0}h)_i\).  On the union of two adjacent
\(t\)-by-\(t\) subcell squares, the ordinary grid Poincare
inequality gives
\begin{equation}\label{eq:op-two-cell}
 |\overline u_{\rm left}-\overline u_{\rm right}|^2
 \leq \frac{2}{t^2}\sum_i|u_i-\overline u|^2
 \leq C\sum_{\langle i,j\rangle}(u_i-u_j)^2.
\end{equation}
The constant does not depend on \(t\).  Apply this to each
same-component observation face and sum with bounded overlap.
The observation residuals then show that
\begin{equation}\label{eq:op-face-control}
 \sum_{\langle x,x'\rangle\ {\rm in}\ G}(y_x-y_{x'})^2
 \leq C_\delta\left\{
       \delta\sum_{\langle i,j\rangle}(u_i-u_j)^2
                          +\norm{B_sh-y}_2^2\right\}.
\end{equation}
Only retained faces within a component occur.

Conversely, \(E_G(y)\) is at most a constant times the left
side of \eqref{eq:op-face-control}.  One explicit trial
interpolates the values \(y_x\) at scale \(s\), with central
plateaus and a piecewise affine subdivision joining cell
centers, face midpoints, and vertices.  At a midpoint or vertex
use the local average within each incident face-connected
sector.  Both its gradient energy and its observation error
are bounded by the nearby face-difference sum.  There is no
bond requirement between opposite-only sectors.  These local
estimates are uniform under the permitted cuts, and agree
under lattice refinement.  They prove the required comparison
of minima.

The constants do not use the microscopic chain precision or
the value of \(\widetilde b\).  Pass to the fixed-graph Skellam
limit using Lemma~\ref{lem:comparisons}.  If necessary decrease
\(c_*\) so that \(c_*\leq1\).  This gives
\eqref{eq:op-large-field} uniformly in the stated parameters.
\end{proof}

\subsection{Entry into the full analytic activity norm}

We use the nearest-neighbor Gaussian split and local subset
algebra of \Bcite{en:entry}.  The theorem there is stated for
quadratic discrete heights; it is \emph{not} being applied
directly to the present physical law.  We prove its needed
entry conclusion using Lemmas~\ref{lem:op-local-limits}
and~\ref{lem:op-large-field}, while retaining its
model-independent Gaussian estimates and finite expansion.

Fix admissible downstream parameters for \Bcite{wh:map} in
the gapped regime \(a'>8\pi\).  Write \(L\) for the block
factor, \(\mathfrak A\) for the activity weight, and \(h_0\)
for the analytic radius.  Choose with strict margins
\begin{equation}\label{eq:op-reserve}
 \max(0,1-c_*)<h_1<h_*<1/p_2,\qquad 1<p_1<p_2.
\end{equation}
The observation-energy regulator and norm are precisely those
of \Bcite{en:variational} and \Bcite{en:norm}, for this fixed
nearest-neighbor reference and these parameters.
For orientation, an activity \(K(X,\psi)\) is a local analytic
correction indexed by a finite set \(X\) of blocks; \(|X|\) is
its block count and \(\psi\) is the reference field.
The point norm \(|K(X)|_{u,X,\psi}\) includes the full sum of
field derivatives at analytic radius \(h_0\). At block side
\(u=mL^j\), the global norm is
\[
 \norm{K}_{u,\mathfrak A,V}
 =\sup_{X,\psi}\mathfrak A^{|X|}
   \frac{|K(X)|_{u,X,\psi}}
        {W_u^\kappa(X,\psi)e^{V_j(X,\psi)}}.
\]
The block supports and ordinary gradient regulator \(W_u^\kappa\)
are those of \Bcite{rg:geometry} and \Bcite{rg:regulator}.
The observation-energy regulator \(V_j\) is displayed in
\eqref{eq:op-regulator} below. Thus the entry estimate controls
all fields and analytic derivatives, not only bounded real data.

\begin{lemma}[Entry for the strengthened law]\label{lem:op-entry}
For every \(\epsilon>0\) there are fixed dyadic integers
\(q,R,t,s_0\), with \(P=R^5\), \(s=ts_0\), and \(m=qRs\),
and a real interval \(I\) about \(b\), such that the exact
relative observation density of the strengthened law, after
averaging the observation noise and reference modes below
scale \(m\), has plane and torus local subset activities
\(K_m\) satisfying
\begin{equation}\label{eq:op-entry-small}
       \norm{K_m(\widetilde b)}_{m,\mathfrak A,V}
                       \leq\epsilon\qquad(\widetilde b\in I).
\end{equation}
The norm includes the full analytic field-derivative sum.
The activities are real, even, square-covariant, invariant
under block-grid translations and common field translations
by \(2\pi/\alpha\).  The torus representation is exact for
compatible sides \(n\geq sP\), with exact plane copying on
unwrapped padded supports and the copying property under
subsequent local maps.
\end{lemma}

\begin{proof}
We give the entry estimates, especially the places where the
physical law differs from the quadratic source.
The choices are made in this order: the downstream norm and
regulator parameters and imaginary strip widths; \(q,R,P\);
a finite large-field threshold; \(t\); \(s_0\) depending on
\(t\); and finally the interval \(I\).

\paragraph{The reference split and exact expansion.}
For this observed law the exact density relative to its
Gaussian reference is, up to normalization,
\begin{equation}\label{eq:op-exact-density}
             Z_{\rm full}(y;\widetilde b)
                         \exp\{E_{\rm full}(y)/2\}.
\end{equation}
This is the observation identity \Bcite{en:exact-density}
with the new \(Z\).  Use the positive finite-range
nearest-neighbor covariance \(C_{<m}\) of
\Bcite{en:entry-covariance}, followed by its positive
finite-range shells.  These covariances, their ranges, and
the reference energies do not involve
\eqref{eq:op-added-precision}.

Use the covariant isolation of \Bcite{en:patches}: tiles of
diameter comparable to \(R\) observation cells, with trimmed
vertex tiles, finite colors of a fixed period \(q\), and
bounded-size proximity groups within each color.  The core,
open patch, and surrounding guard have fixed multiples of
\(R\) as their widths.  At an event, delete all remaining
bonds connecting its tiles to other tiles, counting each
bond once and deleting the corresponding averaged bonds.
If \(o,c\) are the current open patches before and after
deletion, use the replacement
\begin{equation}\label{eq:op-proxy}
 {\cal B}_i(y)=C_i\exp\{-(E_o(y)-E_c(y))/2\},\qquad
       C_i=\frac{Z_o(0;b)}{Z_c(0;b)}.
\end{equation}
The scalar is fixed at \(b\), while the Gaussian shape
always has reference \(a'\).  Replace the old weight by
this proxy times the cut weight plus their difference.
A selected difference marks the event and skips all later
events within a fixed sufficiently large halo.  This is an
exact finite expansion.  A later considered patch avoids
each earlier marked core and guard, so its proxy is
independent of the sign chosen inside that difference.

Divide by the product of the no-mark proxies and centered
isolated-tile normalizers.  Outside mark halos expand the
isolated error
\[
     \frac{Z_Q(y;\widetilde b)e^{E_Q(y)/2}}{Z_Q(0;b)}-1.
\]
The residual Gaussian quadratic difference between the full
energy and the no-mark energy is expanded into links, as
in \Bcite{en:link-matrix} and \Bcite{en:links}.  Gaussian
localization gives link weights \(Ce^{-c(R+d)}\) for length
\(d\), with their full displacement retained on tori.
All averaged-bond supports lie inside the required collars;
they change the physical positive topologies, not these
Gaussian link estimates.

\paragraph{Local normalization and cutting.}
For a marked core, its centered restoration ratio is
sandwiched between the open-patch ratio and the ratio with
its outer guard pinned to zero.  The logarithmic gap is
bounded by \(D(H,p)\), with \(H\) containing every restored
support.  This follows from Lemma~\ref{lem:restoration}
also for the additional positive averaged form.
Lemma~\ref{lem:op-local-limits} and Gaussian localization
make this gap exponentially small in \(R\), after the two
refinements at the fixed geometry.

Only \(O(j_1)\) centered comparisons are charged for
\(j_1\) marks.  Indeed telescope separately the positive
branch and no-mark branch on an enlarged union of their
halos.  The two paths agree at its artificial boundary,
so clipped proxies there cancel; final integrals and proxies
outside the affected union cancel as well.  All events
remaining inside the union number \(O(j_1)\).  This is the
finite double telescoping of \Bcite{en:normalization};
it applies unchanged to the present centered sandwich.
Finite-patch continuity preserves the bounded comparisons
near \(b\), with the constants \(C_i\) fixed as in
\eqref{eq:op-proxy}.  No error proportional to the
unaffected volume is introduced.

For a positive topology \(G\), the resulting bound has
the form
\begin{equation}\label{eq:op-pattern-before-cut}
 e^{Cj_1+Ce^{-cR}{\cal E}(y)}
 e^{E_G(y)/2}\frac{Z_G(y;\widetilde b)}{Z_G(0;\widetilde b)}
                  \E_{G,y}\prod_i\delta_i,
\end{equation}
where \({\cal E}\) is a face-difference cost in the padded
halo and, at prescribed primary guard data \(\gamma_i\),
\[
 \delta_i=
 \frac{|K_o(y,\gamma_i)-{\cal B}_i(y)K_c(y,\gamma_i)|}
      {K_o(y,\gamma_i)+{\cal B}_i(y)K_c(y,\gamma_i)}
                         \leq1.
\]
All factors exterior to a full guard cancel in this ratio.

Cut into \(P\)-boxes on an optimized translate of the
\(R\)-grid.  Averaging translates loses at most
\(C(R/P)j_1\) marks and costs at most
\(C(R/P){\cal E}(y)\) in Gaussian energy.  The joint
shifted-versus-centered comparison at prescribed guard
data is Lemma~\ref{lem:comparisons}, with the observations
embedded as height-plus-noise coordinates.  Its remaining
centered guard factors cost \(e^{Cj_1}\): impose the
guards sequentially and use open and zero-guard bounds
with separated outer collars.  Thus the cut comparison
has the same seam charge as \Bcite{en:seam-cost},
\begin{equation}\label{eq:op-seam}
                        e^{Cj_1+C(R/P){\cal E}(y)}.
\end{equation}

\paragraph{Bounded data and large data.}
We distinguish the graph used to classify observation data
from the positive topologies being compared.  Let \(D^0\)
be the observation-cell graph on the pattern's original padded
support \(D\), \emph{before any isolation deletion}.  Clip this
graph by the \(P\)-box grid and put
\begin{equation}\label{eq:op-original-face-cost}
 {\cal E}_{\rm box}(y)
 =\sum_{\substack{\langle x,x'\rangle\text{ in }D^0\\
                         x,x'\text{ in the box}}}
                                  (y_x-y_{x'})^2.
\end{equation}
In particular, a face remains in this cost when an isolation
step deletes its physical bonds.  Retain a deep mark only if
its complete open patch, restored supports, and guard lie in
the same clipped box; the other marks are among the boundary
losses already charged in \eqref{eq:op-seam}.  Thus every face
across which any retained restoration, Gaussian proxy, or
guard comparison tests a relative translation belongs to this
original graph.  Clipping adds no observation cell or exterior
field value to \(D\).

Fix the large-field threshold temporarily.  On each connected
component of this \emph{original} clipped graph translate all
heights, observations, and guard data together by one height
period, putting one observed value in \([-\pi,\pi]\).  Every
factor under comparison is invariant under this common
translation.  Bounded \({\cal E}_{\rm box}\) now gives a
compact set of observation values by following paths in the
finite original graph.  Independent translations of components
created by a later cut are not used: a restoration may join
those components and detect their relative translation.  At
fixed \(R,P\) there are finitely many geometries, positive
topologies, and mark choices.  Each Gaussian energy and each
physical partition function still uses its own actual positive
topology, not \(D^0\).

Decompose the logarithm of
\(K_o/({\cal B}_iK_c)\) into its centered zero-guard value,
the variation of its centered restoration with guard data,
and its two conditional observation Laplace terms minus
the Gaussian proxy change.  The first term is
\(O(e^{-cR})\) after refinement.  The second lies between
\(\log r_H\) and zero.  The conditional terms and \(r_H\)
converge by Lemma~\ref{lem:op-local-limits} under either
topology and its fixed observation tilt.

In the limiting Gaussian energy, let \(u_c\) be the cut
harmonic prediction from guard data and let \(\chi_H\)
equal one at the core and vanish before the guard.
The unchanged nearest-neighbor localization estimate
\Bcite{en:prediction} gives, for every \(r\geq2\),
\begin{equation}\label{eq:op-prediction}
 \left\|\norm{\chi_Hu_c}_T\right\|_{L^r}
 \leq C\sqrt r\,{\rm poly}(P)e^{-cR}
                         (1+\sqrt{{\cal E}_{\rm box}(y)}).
\end{equation}
The squared cutoff energy bounds \(-2\log r_H\).
Correcting a cut minimizer by this cutoff makes it
admissible for restoration without changing the guard;
orthogonal projection therefore bounds the linear
conditional-Laplace discrepancy by
\(C\norm{y}_2\norm{\chi_Hu_c}_T\).
The conditional quadratic discrepancy from the open-patch
proxy is bounded by \(Ce^{-cR}\norm{y}_2^2\), by the same
guarded localization.  These arguments use finite-energy
predictions and finite guard probes, not pointwise
continuum guard values.

Transfer these estimates using the conditional approximation
and \(r_H\leq r_\Delta\leq1\).  The Gaussian limiting
\(r_H\) is strictly positive, so exceptional sets are
negligible in both topologies, also after fixed tilts.
On such sets use \(\delta_i\leq1\).  Since
\(|(x-1)/(x+1)|\leq\min(1,|\log x|)\), bounded moments
of each \(\delta_i\) have the bounds from
\eqref{eq:op-prediction}.
Holder over the at most \({\rm poly}(P)\) deep marks,
and absorption of polynomial losses, give
\begin{equation}\label{eq:op-good-box}
 C\exp\{-cRj_{\rm deep}
                       +R^{-12}{\cal E}_{\rm box}(y)\}.
\end{equation}
For example, maximize
\((1+{\cal E})^{Cj}e^{-R^{-12}{\cal E}}\).
Since \(j\leq{\rm poly}(P)\), a factor
\([P^Cj^C(1+{\cal E})^C]^j\) is bounded by
\(e^{O(j\log R)+R^{-12}{\cal E}}\).
The \(O(j\log R)\) loss is absorbed by the exponential
rarity.  We use this smaller allowance so that the conservative
\(P^2\) coordinate count is negligible below.
Isolated selected-tile errors
can be made smaller than \(e^{-cR}\) by the same compact
Gaussian convergence.

We also need a separate estimate after a difference is expanded
into its two positive terms.  For every actual positive box
topology \(G\), Lemma~\ref{lem:op-local-limits} gives
\[
 e^{E_G(y)/2}\frac{Z_G(y;b)}{Z_G(0;b)}\longrightarrow1
\]
uniformly on the good-data compact set just described, after
the common period recentering.  The reference limit is precisely
the Gaussian of coefficient \(a'\), on that same topology.
There are only finitely many such topologies.  Choose the
refinements so that, on every good box,
\begin{equation}\label{eq:op-positive-good-box}
 e^{E_G(y)/2}\frac{Z_G(y;\widetilde b)}{Z_G(0;\widetilde b)}
       \leq 2\exp\{R^{-12}{\cal E}_{\rm box}(y)\}.
\end{equation}
At \(\widetilde b=b\) choose the stronger upper bound \(3/2\)
using that convergence; finite-patch continuity gives the bound \(2\)
after shrinking the eventual parameter interval.  For an
isolated-tile error expanded into its positive terms, the
constant Gaussian term satisfies the same bound.  Centered
normalizations between branches cost the already proved
\(e^{Cj_1}\).  Equation~\eqref{eq:op-positive-good-box}
therefore bounds positive branches without any discrepancy
factor \(\delta_i\); it is not obtained by removing
\(\delta_i\) from \eqref{eq:op-good-box}.  The displayed
allowance is retained to use the same quadratic-cost accounting
in both estimates.

On a bad box discard rarity and apply
Lemma~\ref{lem:op-large-field}.  Its centered-normalized
factor is at most
\[
 e^{E_G(y)/2}\frac{Z_G(y;\widetilde b)}{Z_G(0;\widetilde b)}
                     \leq e^{(1-c_*)E_G(y)/2}.
\]
This replaces the bare quadratic estimate
\Bcite{en:bad-box}.  It is the only large-field input
needed from the physical height law.

\paragraph{Gaussian averaging with reserve.}
Here is the energy accounting for the original-face costs.
For any real trial field \(v\) on the uncut padded domain,
two-cell Poincare and the observation residual give
\begin{equation}\label{eq:op-original-face-control}
 \sum_{\langle x,x'\rangle\text{ in }D^0}(y_x-y_{x'})^2
 \leq C\{a'^{-1}\ip{v}{A_Dv}+\norm{B_sv-y}_2^2\}.
\end{equation}
Indeed insert the two cell averages of \(v\) in each
difference.  The difference of averages is bounded by the
gradient energy on the two adjacent cells, as in
\eqref{eq:op-two-cell}; the two residuals are charged at
most four times.  All faces are internal to \(D^0\), so
these two-cell regions are present before minimization and
have bounded overlap, also at a clipped boundary.  Minimizing
over \(v\) proves the bound by \(CE_D(y)\).  This argument
does not ask a cut topology to control a deleted face.

For every positive branch of the expansion, restriction of
that same uncut trial field to the disjoint cut boxes bounds
the sum of their actual minimized energies \(E_G\) by
\(E_D\).  The bad-box physical contribution is consequently
at most \((1-c_*)E_D/2\).  All remaining data costs---the
\(R^{-12}{\cal E}_{\rm box}\) good-box allowance, seam
and halo costs, Gaussian links, and
\(\varepsilon_1{\cal E}_{\rm box}\) bad-box threshold
charges---use faces of \(D^0\).  Their summed coefficients
are at most
\(C(R^{-12}+R/P+e^{-cR}+\varepsilon_1)\) per face.
For halos this follows from their bounded multiplicity;
for links one sums their exponentially decaying displacement
weights.  The same bound applies to every restoration branch.
Choose \(R\) large and then \(\varepsilon_1>0\) small so
that \eqref{eq:op-original-face-control} makes their total
less than \([h_1-(1-c_*)]E_D/2\).  Thus the total positive
quadratic cost is at most \(h_1E_D/2\), with strict reserve
as required by \eqref{eq:op-reserve}.  The large-field
threshold is chosen only after this reserve is fixed.

The reference covariance contracts the observation
energy, exactly as in \Bcite{en:energy-contraction}.
At scale \(u=mL^j\), the regulator is
\begin{equation}\label{eq:op-regulator}
 V_j(X,\psi)=\frac12\sup_{\zeta,e}
 \left\{h_*E_{D_u(X)}\bigl(B_s\alpha(\psi+\zeta)+e\bigr)
       -p_1^{-1}\bigl(\norm{\zeta}_{<u}^2+\norm e_2^2\bigr)\right\}.
\end{equation}
Here the first norm is the Cameron norm of the reference
pieces below \(u\); singular covariances mean their
Gaussian Hilbert spaces.  The covariance contraction
and \(p_1h_*<1\) make this supremum finite.  Its
composition, localization, and terminal bounds are
\Bcite{en:regulator}; these involve only the unchanged
reference covariance and whole observation cells.

The determinant actually integrated is charged by its
localized quadratic costs, not by the larger dominating
form \(h_1E_D\).  Its trace is bounded by
\begin{equation}\label{eq:op-trace}
 C\,{\rm poly}(P)N_{\rm bad}
 +CR^{-12}P^2N_{\rm good}
 +C(R/P)R^2(j_1+j_2)
 +C\sum_\ell e^{-c(R+d_\ell)}{\rm poly}(d_\ell).
\end{equation}
Here \(N_{\rm good}\) counts occupied good boxes.
Broad bad-box terms, including their threshold charges,
have rank \(O(P^2)\).  The good-box allowance in
\eqref{eq:op-good-box} contributes at most
\(CR^{-12}P^2\) per occupied box, counting all present
cells of the clipped box.  Here \(P^2\) is only an upper
bound on their number; the actual quadratic cost remains
inside \(D\).  A halo intersects a bounded number of
\(P\)-boxes, so \(N_{\rm good}\leq C(j_1+j_2)\).
The remaining non-broad halo terms have coefficient
\(O(R/P)\) on \(O(R^2)\) observation cells, and links
retain their exponentially small path costs.
The added physical precision is contained inside \(Z_G\);
it is not a term of this reference Gaussian cost matrix.
Covariance contraction
gives a spectral margin below one, so logarithmic
determinant costs are bounded by a fixed multiple of
this trace.  No determinant counts subcells or
microscopic height variables.

Insert on a bad box the upper bound
\[
 1\leq\exp\{\varepsilon_1(
                    {\cal E}_{\rm box}(y)-U_{\rm big}^2)\},
\]
with sufficiently small fixed \(\varepsilon_1>0\).
Choose the finite \(U_{\rm big}\) to pay the polynomial
bad-box losses, the first term of \eqref{eq:op-trace},
and the lost item gains.  This choice precedes \(t,s_0\).
Since \(P=R^5\), both the full good-box allowance
\(R^{-12}P^2\) and the halo seam trace cost are
\(O(R^{-2})\) per item.  Gaussian completion with reserve consequently
gives, for an averaged pattern \(F_\pi\),
\begin{equation}\label{eq:op-real-pattern}
 |F_\pi(\psi)|
 \leq C^{j_1+j_2}e^{-cR(j_1+j_2)}
       \prod_\ell Ce^{-c(R+d_\ell)}
                     W_m^{\kappa_0}(X,\psi)e^{V_0(X,\psi)},
\end{equation}
where \(\kappa_0<\kappa\) leaves ordinary regulator
reserve.  A real translate by \(xf\), with \(f\) of
scaled point norm at most one, costs at most
\(e^{Cx^2|X|}\), with the final regulator still evaluated
at \(\psi\).  This uses \(h_1<h_*\) in the optimized
quadratic form and
\(\ip{f}{A_Df}\leq C|X|\).

\paragraph{Imaginary directions and analytic derivatives.}
The full analytic norm cannot be inferred from
bounded-real-data convergence alone.  On the extended
quadratic lattice, including
\eqref{eq:op-added-precision}, the theta parity identity
\Bcite{en:theta-parity} gives
\begin{equation}\label{eq:op-complex-comparison}
 \frac{|Z_G(y+iz;\widetilde b)|}{Z_G(y;\widetilde b)}
             \leq\frac{Z_G(iz;\widetilde b)}{Z_G(0;\widetilde b)}.
\end{equation}
It survives the fixed-graph Skellam limit.  The centered
imaginary ratio is its Gaussian prefactor times a
positive characteristic function; centered pinning
increases that characteristic factor.

Expand differences into positive topologies, temporarily
discarding their rarity, and apply
\eqref{eq:op-complex-comparison} to each whole integrated
topology.  Pin strips of width \(R/4\) along translated
\(P\)-box walls.  All walls and strip boundaries are
on observation-cell boundaries: \(R\) is dyadic,
\(R/4\) is integral, and the translates use the
\(R\)-grid.  Thus strips are unions of whole
observation cells and of their nested subcells;
no observation average is split.  Cutting through their interiors
disconnects also the averaged bonds, because all their
supporting subcells there are pinned.  On the resulting
bounded boxes, Lemma~\ref{lem:op-local-limits} gives the
Gaussian imaginary ratios uniformly on bounded
imaginary directions, with an arbitrarily small
multiplicative error after refinement.  Their limits
are strictly positive on each such compact set.

The additional Gaussian wall energy, averaged over wall
translates, is at most
\[
             C\norm z_\infty^2(R/P)
                         \#\{\hbox{active observation cells}\},
\]
by \Bcite{en:imaginary-wall}.  The other imaginary
Gaussian shapes cancel those of the proxies.  Real
parts are bounded on good boxes by the separate
positive-topology estimate \eqref{eq:op-positive-good-box},
and on bad boxes by Lemma~\ref{lem:op-large-field}.
Apply the preceding Gaussian completion to these bounds,
including the explicitly charged
\(CR^{-12}P^2N_{\rm good}\) term in
\eqref{eq:op-trace}; links retain their gains.
For fixed strip height \(H\) we obtain
\begin{align}\label{eq:op-rough-strip}
 |F_\pi(\psi+(x+iy)f)|
 &\leq
 \exp\{o(R)(j_1+j_2)+C_H(1+x^2)|X|\}
       \prod_\ell Ce^{-c(R+d_\ell)}
       W_m^\kappa(X,\psi)e^{V_0(X,\psi)},\nonumber\\[-2pt]
 &\hspace{65mm}|y|\leq H .
\end{align}
In particular, the wall charge per halo is
\(O(R^3/P)=O(R^{-2})\), not \(O(R^2)\).
Absolute Gaussian integration and analyticity follow
from \eqref{eq:op-large-field} with reserve and these
strip bounds.

Choose \(r_0>4eh_0\), then \(H>4r_0\), before \(R\).
Apply three-lines on the upper and lower strips to
\(F_\pi(\psi+zf)e^{-D_0z^2|X|}\), with \(D_0>C_H\).
The real boundary has \eqref{eq:op-real-pattern} and
the real-translation bound; the other boundary has
\eqref{eq:op-rough-strip}.  On \(|z|\leq r_0\) the
real-boundary harmonic weight is at least \(3/4\), so
the exponential item gains survive with a smaller
constant.  Removing the damping costs
\(e^{C|X|}\), absorbed because
\[
 |X|\leq C(j_1+j_2)+C\sum_\ell(1+d_\ell/R).
\]
Cauchy's estimate and real polarization bound the
order-\(r\) contribution to the analytic point norm
by the pattern bound times
\((h_0/r_0)^r r^r/r!\).  Its sum converges.  Thus
\eqref{eq:op-real-pattern}, with weaker constants,
holds in the full analytic point norm.

\paragraph{Uniformity and completion of entry.}
At fixed \(R,P\), choose a sufficiently large \(t\),
then a sufficiently large \(s_0\) depending on \(t\),
using Lemma~\ref{lem:op-local-limits} on the finitely
many bounded geometries, real-data compact sets, and
imaginary-direction compact sets just used.  All
estimates have strict slack.  Fixed-graph sums and
conditional bounded integrals are continuous in
\(\widetilde b\).  Indeed guard dimension is now
finite and fixed.  Although guard values range over
an unbounded lattice, \(\delta_i\leq1\), and the joint
laws and fixed tilts have uniform spanning-forest
exponential domination on the compact parameter and
observation sets.  Their positive normalizing sums
have positive compact lower bounds.  Dominated
convergence therefore applies to the integrated
absolute differences; no pointwise uniform estimate
at arbitrary guard data is needed.  Hence a
common real neighborhood of \(b\) preserves these
bounded-data estimates.  The global large-field
bound already has a uniform reserve.  The same
Gaussian averaging and interpolation therefore give
the full norm bound throughout this neighborhood;
no continuity assertion about the unknown \(a(\widetilde b)\)
has entered.

Finally use the connected pattern summation of
\Bcite{en:placement}.  A halo has bounded block
diameter and only polynomially many observation-cell
labels per block.  A link of length \(d\) covers
at most \(C(1+d/R)\) blocks.  The exponential item
and link gains pay all activity weights, coverings,
and rooted-tree sums, leaving \(Ce^{-cR}\) in the
norm.  Choose \(R\) large enough for
\eqref{eq:op-entry-small}.  Absolute estimates
justify every exchange with Gaussian integration.
Finite-range covariance factorization gives the
exact subset algebra, including its exclusions.
The covariant tile rules and aligned added form
preserve the stated symmetries.  The activity hypotheses
preceding \Bcite{wh:map} require block translations
and square symmetries, not unit-site translation
invariance of the physical interaction; the nested
grids have \(s_0\mid s\mid m\) and common square
centers.  Winding links keep
their full covers; every unwrapped padded support
therefore has exactly its plane activity.  This
proves the lemma.
\end{proof}

\subsection{Uniform gapped stability and a torus diagnostic}

We finish the proof of Proposition~\ref{prop:openness}.
Apply \Bcite{wh:map} with the cosine coordinate omitted,
which is permitted because \(a'>8\pi\).  The reference
kernel is fixed and nearest-neighbor; take its starting
scale sufficiently large.  The modified regulator is
allowed by \Bcite{en:regulator}.  All map hypotheses now
concern the reference and the small activities supplied
by Lemma~\ref{lem:op-entry}.

We record why the permitted smallness can be chosen
\emph{before} the large entry scale.  The Gaussian-curve
construction \Bcite{ep:gaussian-curves} gives exact
trajectories by initializing a small multiple \(d\)
of the block gradient energy and no lattice factor.
Their expansion has a linear gradient singleton and
an \(O(d^2)\) remainder.  These estimates are uniform
over sufficiently large starting scales: scaled block
gradient energy has a uniform point norm, and its
small exponential is absorbed by the same regulator
reserve.  Terminal Gaussian completion is also uniform,
because the tail covariance contracts the kinetic
energy and the terminal torus has a fixed number of
blocks.  On a fixed complex disk,
\((\Id-dC^{1/2}AC^{1/2})^{-1}\) is uniformly bounded,
and determinants cancel in the shifted ratios.
The map contraction and terminal gradient diagnostic
then prevent a first exit on that disk, exactly as
in the proof of \Bcite{ep:gaussian-curves}.  Its
analytic inverse disk is consequently uniform in
the starting scale as well as the RG step.  This
uniformity is a consequence of that proof and
\Bcite{en:regulator}, rather than an extra physical
assumption on the strengthened model.

Use the exact Gaussian parameter \(u\) for the
gradient coordinate and subtract the corresponding
Gaussian remainder, writing the transverse remainder
as \(B_{\rm rem}\).  The balanced linearization is
contracting transversely and the pure Gaussian curve
has no drift.  The analytic map estimates give
\begin{align}
 \norm{B_{{\rm rem},+}}
 &\leq[\theta+C(|u|+\norm{B_{\rm rem}})]
                                      \norm{B_{\rm rem}},\label{eq:op-gapped-B}\\
 |u_+-u|
 &\leq C(|u|+\norm{B_{\rm rem}})
                                      \norm{B_{\rm rem}}.\label{eq:op-gapped-u}
\end{align}
These are the fixed-reference inequalities
\Bcite{ep:gapped-B} and \Bcite{ep:gapped-u}.
Choose a small tube so that the first factor is at
most \(\theta'<1\).  Small enough entry data stay
in that tube, with
\(\norm{B_{{\rm rem},j}}\leq
(\theta')^j\norm{B_{{\rm rem},0}}\) and summable
gradient drift.  In particular all activities
remain as small as desired.  The separate torus
remainder obeys the same uniform small bound by
its contraction with quadratic forcing and the
exact copying of singleton coefficients in
\Bcite{wh:map}.

Choose a sufficiently large fixed dyadic \(D\) for
the causal neighborhoods of these maps and use
compatible tori of side
\begin{equation}\label{eq:op-stopping-tori}
                        n_j=DmL^j .
\end{equation}
Take \(j\) large enough that \(n_j\geq sP\), so
the exact entry is available.  At the stopping
scale there are \(D^2\) blocks.  Let \(f\) be a
nonzero mean-free first cosine on the unit torus,
and let \(f_n\) be its area-weighted observation-cell
samples.  Put
\[
                   g=\ip{f}{(-\Delta_{\T^2})^{-1}f}>0.
\]
The reference variance of \(\ip{Y}{f_n}\) tends
to \(a'g\).  Gaussian completion of a smooth real
source translates the exact relative density by
the reference covariance of that source.  The
Gaussian Laplace prefactor remains, while the
fixed terminal block count and regulator reserve
bound the differentiated logarithmic correction
by the small terminal activity norms.  This is
the gradient diagnostic
\Bcite{ep:gradient-diagnostic}.  It requires only
smallness here, not convergence of the gradient
coordinate to zero and not a derivative in
\(\widetilde b\).  By choosing the tube small,
for all sufficiently large \(j\),
\begin{equation}\label{eq:op-positive-torus}
 \Var_{\rm strengthened}\ip{Y}{f_{n_j}}
                                      \geq \tfrac12 a'g .
\end{equation}
The observation noise has variance
\(\sum_x f_{n_j}(x)^2=O((s/n_j)^2)\), which tends
to zero.  Removing the added positive precision
increases centered height variance by
Lemma~\ref{lem:comparisons}.  Hence the original
periodic Skellam height law satisfies
\begin{equation}\label{eq:op-original-torus}
 \liminf_{j\to\infty}
 \Var_{\widetilde b,{\rm tor}}
             \ip{h}{B_s^*f_{n_j}}>0
                    \qquad(\widetilde b\in I).
\end{equation}

If \(a(\widetilde b)=0\), the left side must instead
vanish.  To see this directly, first take a smooth
periodic gradient test.  Split the torus into
macroscopic rectangles, release the bonds between
them, and retain the internal gradient observations.
For a fixed macroscopic partition there are \(O(n)\)
seam edges and their coefficients are \(O(n^{-1})\),
so the discarded seams have reference cost \(O(n^{-1})\);
the free-cell variances converge to zero by
Theorem~\ref{thm:free-limits}.  The variance comparison
and Lemma~\ref{lem:reference-cost} give vanishing
variance for the original periodic test.  The
same conclusion holds for smooth mean-zero
density tests by using the periodic inverse
Laplacian and reference-energy approximation.
Finally, the observation-cell profile
\(B_s^*f_n\) approximates the corresponding
smooth density in reference cost as \(n\to\infty\),
with reference-cost error \(O((s/n)^2)\),
since \(s\) is fixed.  Equivalently one may use
the sampled periodic cosine shift and estimate
its cell-averaging error.  Thus its variance
also tends to zero, contradicting
\eqref{eq:op-original-torus}.

It follows that \(a(\widetilde b)>0\) throughout
the real neighborhood \(I\).  Shrinking \(I\)
inside \((0,\infty)\) proves
Proposition~\ref{prop:openness}.
